\documentclass[12pt, a4paper]{article}

% Paquetes requeridos para el circuito
\usepackage[utf8]{inputenc}
\usepackage[T1]{fontenc}
\usepackage{circuitikz} 

\begin{document}

\begin{center}
    \begin{circuitikz}[american, scale=1.0, transform shape]
        
        % 1. Rama 1 (Izquierda): Fuente V y Resistencia R1 en serie desde 'b' a 'a'
        \draw (0,0) to[V, l=$V$] (0,2)
                    to[R, l=$R_1$] (0,4);
        
        % 2. Rama 2 (Centro): Fuente de corriente I en paralelo
        \draw (2.5,0) to[I, l=$I$] (2.5,4);
        
        % 3. Rama 3 (Derecha): Resistencia R2 en paralelo
        \draw (5,0) to[R, l=$R_2$] (5,4);
        
        % Conexiones horizontales para los nodos comunes superior e inferior
        \draw (0,4) -- (5,4);
        \draw (0,0) -- (5,0);
        
        % 4. Terminales explícitos de salida para 'a' y 'b'
        \draw (5,4) to[short, -o] (6.5,4) node[right] {a};
        \draw (5,0) to[short, -o] (6.5,0) node[right] {b};
        
    \end{circuitikz}
\end{center}

\end{document}